% ============================================================
%  LAMPIRAN 2 - VISUALISASI KNOWLEDGE GRAPH PER MATA PELAJARAN
%  Compile standalone:  pdflatex lampiran-2.tex
%  Compile as part of thesis: pdflatex main.tex
% ============================================================
\documentclass[../../thesis]{subfiles}
\usepackage{hyphenat}
\usepackage{iftex}

% ── Encoding, Language & Fonts ───────────────────────────────
\ifXeTeX
    \usepackage{polyglossia}
    \setmainlanguage{bahasai}
    \usepackage{fontspec}
    \setmainfont{TeX Gyre Termes}
    \setsansfont{TeX Gyre Heros}
    \setmonofont[Scale=0.88]{TeX Gyre Cursor}
\else
    \usepackage[utf8]{inputenc}
    \usepackage[T1]{fontenc}
    \usepackage[indonesian]{babel}
    \usepackage{mathptmx}
    \usepackage{helvet}
    \usepackage{courier}
\fi

% ── Page geometry (A4, UI thesis margins) ────────────────────
\usepackage[
  a4paper,
  top    = 2.54cm,
  bottom = 2.54cm,
  left   = 3.17cm,
  right  = 2.54cm
]{geometry}

% ── Line spacing & paragraph ─────────────────────────────────
\usepackage{setspace}
\onehalfspacing
\setlength{\parindent}{1.27cm}
\setlength{\parskip}{0pt}

% ── Microtype & justification ────────────────────────────────
\usepackage{microtype}
\sloppy

% ── Headings ─────────────────────────────────────────────────
\usepackage{titlesec}

\titleformat{\chapter}[block]
  {\centering\rmfamily\bfseries\large}{}
  {0pt}{\MakeUppercase}
\titlespacing*{\chapter}{0pt}{0pt}{1.5em}

\titleformat{\section}
  {\rmfamily\bfseries\normalsize}{\thesection}{1em}{}
\titlespacing*{\section}{0pt}{1.2em}{0.6em}

% ── Graphics ─────────────────────────────────────────────────
\usepackage{graphicx}
\makeatletter
\def\maxwidth{\ifdim\Gin@nat@width>\linewidth\linewidth\else\Gin@nat@width\fi}
\makeatother
\setkeys{Gin}{width=\maxwidth,keepaspectratio}

% ── Hyperlinks ───────────────────────────────────────────────
\usepackage[hidelinks]{hyperref}

% ============================================================
\begin{document}

% When standalone: switch to appendix mode and set the counter
% so appendix-local numbering renders correctly. When compiled inside main.tex,
% main.tex issues \appendix and advances the counter before \subfile.
\ifSubfilesClassLoaded{}{\appendix\setcounter{chapter}{2}}

\lampiran{Visualisasi Knowledge Graph per Mata Pelajaran}
\label{lamp:kg-per-mapel}

Lampiran ini menyajikan visualisasi topologi \textit{knowledge graph}
hasil konstruksi untuk masing-masing mata pelajaran (Fisika, Kimia, dan
Biologi) pada kondisi sebelum proses \textit{Knowledge Graph Completion}
(intra-mapel). Visualisasi ini melengkapi analisis struktural pada
Subbab~\ref{subsec:struktural} dengan memperlihatkan struktur internal
tiap disiplin secara terpisah.

\begin{figure}[htbp]
  \centering
  \includegraphics[width=\linewidth]{assets/pics/kg/fisika.png}
  \caption{Topologi \textit{knowledge graph} mata pelajaran Fisika.}
  \label{fig:kg-fisika}
\end{figure}

\begin{figure}[htbp]
  \centering
  \includegraphics[width=\linewidth]{assets/pics/kg/kimia.png}
  \caption{Topologi \textit{knowledge graph} mata pelajaran Kimia.}
  \label{fig:kg-kimia}
\end{figure}

\begin{figure}[htbp]
  \centering
  \includegraphics[width=\linewidth]{assets/pics/kg/biologi.png}
  \caption{Topologi \textit{knowledge graph} mata pelajaran Biologi.}
  \label{fig:kg-biologi}
\end{figure}

\end{document}
